\clearpage
\begin{proofbox}
\textbf{证明一（优先）：代入一次，展开一次，再用韦达。}
若 $B\ne0$，由直线得到 $y=-(Ax+C)/B$。代入并乘以正数 $\alpha\beta B^2$：
\[
 \frac{x^2}{\alpha}+\frac{(Ax+C)^2}{\beta B^2}=1
 \ \Longleftrightarrow\
 \beta B^2x^2+\alpha(Ax+C)^2=\alpha\beta B^2.
\]
展开即为
\[
 (A^2\alpha+B^2\beta)x^2+2AC\alpha x+\alpha C^2-\alpha\beta B^2=0.
\]
二次项系数 $D>0$，由韦达定理立即得到结论中的横坐标和、积。
反过来，任一实根 $x$ 都唯一对应 $y=-(Ax+C)/B$，满足原来的两条方程；没有增根或漏点。
若 $A\ne0$，同样代入 $x=-(By+C)/A$，得到
\[
 \alpha A^2y^2+\beta(By+C)^2=\alpha\beta A^2,
\]
展开即得纵坐标方程与韦达关系。

\textbf{判别式也可化成一行。} 对满足 $B\ne0$ 的横坐标方程，
\[
\begin{aligned}
 \Delta_x
 &=(2AC\alpha)^2-4D\alpha(C^2-B^2\beta)\\
 &=4\alpha\bigl[C^2(A^2\alpha-D)+DB^2\beta\bigr]\\
 &=4\alpha\beta B^2(D-C^2)=4\alpha\beta B^2H.
\end{aligned}
\]
因为前面的因子严格为正，判别式与 $H$ 同号。
对 $A\ne0$ 的纵坐标方程同理有 $\Delta_y=4\alpha\beta A^2H$。
$A,B$ 至少一个非零，总能选合法消元方向，所以 $H$ 的三分判定覆盖所有直线。

\textbf{韦达本身也很直接：} 若 $Dt^2+qt+r=D(t-t_1)(t-t_2)$，比较一次项与常数项，
便有 $t_1+t_2=-q/D$、$t_1t_2=r/D$。

\tcbbreak
\textbf{第二个角度（几何）：缩放为单位圆，直接看交点条件。}
\begin{center}\begin{tikzpicture}[font=\small]
\draw[black!30,rounded corners=2pt] (0,0) rectangle (14.4,8.2);
\node[anchor=west,font=\small\bfseries] at (.5,7.72) {第二个角度：把椭圆缩放为单位圆，再看直线与圆的位置};
% MODEL 9 4 0 1 -1
\begin{scope}[shift={(4,3.6)},scale=0.85]
\draw[blue,very thick] (0,0) ellipse (3.0000000000 and 2.0000000000);
\draw[orange!85!black,very thick] (-2.5980762114,1)--(2.5980762114,1);\coordinate (M) at (0.0000000000,1.0000000000);
\fill[black] (M) circle (1.8pt) node[above=6pt] {$M$};

\end{scope}
\begin{scope}[shift={(10.6,3.6)},scale=1.55]
\draw[blue,very thick] (0,0) circle (1);\draw[orange!85!black,very thick] (-.8660254038,.5)--(.8660254038,.5);\draw[green!45!black,dashed] (0,0)--(0,.5);\coordinate (N) at (0.0000000000,0.5000000000);
\coordinate (O) at (0.0000000000,0.0000000000);
\fill[black] (N) circle (1.8pt) node[above=6pt] {$M'$};
\fill[black] (O) circle (1.8pt) node[below=6pt] {$O$};

\end{scope}
\node at (4,6.6) {原椭圆：$x^2/9+y^2/4=1$};\node at (10.6,6.6) {变换后：$u^2+v^2=1$};\draw[-Stealth,thick] (7,3.6)--(8.3,3.6);\node[anchor=west] at (.8,.6) {$u=x/3,\ v=y/2$：横纵缩放不同；右侧显示单位圆};
\end{tikzpicture}\end{center}

令 $u=x/\sqrt\alpha,v=y/\sqrt\beta$，椭圆与直线变成
\[
 u^2+v^2=1,\qquad (A\sqrt\alpha)u+(B\sqrt\beta)v+C=0.
\]
变换可逆，不改变交点个数。新直线到原点的距离为
\[
 d=\frac{|C|}{\sqrt{A^2\alpha+B^2\beta}}=\frac{|C|}{\sqrt D}.
\]
单位圆中，$d<1,d=1,d>1$ 分别为割线、切线、无交点；即 $H>0,H=0,H<0$。
当 $H>0$ 时，圆心到弦的垂足就是弦中点；$H=0$ 时该点为切点。其新坐标为
\[
 M'=\left(-\frac{AC\sqrt\alpha}{D},-\frac{BC\sqrt\beta}{D}\right).
\]
例如令法向量为 $(A\sqrt\alpha,B\sqrt\beta)$，垂足在该方向上，代入直线就得到此式。
缩放保留中点，因此变回原坐标得到
\[
 M=\left(-\frac{AC\alpha}{D},-\frac{BC\beta}{D}\right).
\]
\textbf{注意：} 非等比例的横纵缩放一般不保角，也不按同一倍数改变所有长度。
原椭圆中 $OM$ 不一定垂直于弦；下一种证明不依赖这个错误假设。

\tcbbreak
\textbf{第三个角度：两点作差，直接找中点；再求弦长。}
设已有两个不同交点。两点各满足椭圆方程，相减并分解平方差，得
\[
 \frac{(x_1-x_2)(x_1+x_2)}{\alpha}
 +\frac{(y_1-y_2)(y_1+y_2)}{\beta}=0.
\]
记中点为 $(x_M,y_M)$。弦的方向向量与 $(B,-A)$ 平行，故可约去一个非零比例因子，得到
\[
 \frac{B x_M}{\alpha}-\frac{A y_M}{\beta}=0,\qquad
 Ax_M+By_M+C=0.
\]
联立解出 $x_M=-AC\alpha/D,y_M=-BC\beta/D$。
这里没有除以 $A$ 或 $B$，水平和竖直弦都适用。

沿直线写任意点 $(x,y)=M+t(B,-A)$。将它代入椭圆：
\[
 \frac{x^2}{\alpha}+\frac{y^2}{\beta}
 =\frac{C^2}{D}+2t\left(\frac{Bx_M}{\alpha}-\frac{Ay_M}{\beta}\right)
 +t^2\left(\frac{B^2}{\alpha}+\frac{A^2}{\beta}\right)
 =\frac{C^2}{D}+\frac{D}{\alpha\beta}t^2.
\]
所以 $t^2=\alpha\beta H/D^2$。$H>0$ 时两端点分别取正负值，方向向量的长度为
$\sqrt{A^2+B^2}$，于是
\[
 \boxed{PQ=\frac{2\sqrt{\alpha\beta(A^2+B^2)H}}{D}}.
\]
这同时说明了实交点条件和水平、竖直方向为何不需要另造弦长公式。

\tcbbreak
\textbf{若已用斜率消元，根差怎样变成弦长？}
\begin{center}\begin{tikzpicture}[font=\small]
\draw[black!30,rounded corners=2pt] (0,0) rectangle (14.4,8.2);
\node[anchor=west,font=\small\bfseries] at (.5,7.72) {根差只是水平距离：用一个直角三角形补回真正弦长};
% MODEL 8 4 1 1 -1
\begin{scope}[shift={(7.2,3.65)},scale=1.08]
\coordinate (P) at (-0.8968052533,1.8968052533);
\coordinate (Q) at (2.2301385866,-1.2301385866);
\coordinate (M) at (0.6666666667,0.3333333333);
\coordinate (O) at (0.0000000000,0.0000000000);
\coordinate (T) at (2.2301385866,1.8968052533);
\draw[blue,very thick] (0,0) ellipse (2.8284271247 and 2.0000000000);
\draw[orange!85!black,very thick] (P)--(Q);\draw[green!45!black,thick,dashed] (P)--(T)--(Q);\pic[draw,angle radius=3mm] {right angle=P--T--Q};\fill[black] (P) circle (1.8pt) node[above left=6pt] {$P$};
\fill[black] (Q) circle (1.8pt) node[below right=6pt] {$Q$};
\fill[black] (T) circle (1.8pt) node[above right=6pt] {$T$};
\path (P)--(T) node[midway,above=8pt] {$|x_2-x_1|$};\path (T)--(Q) node[midway,right=28pt] {$|y_2-y_1|$};
\end{scope}
\node[anchor=west] at (.8,.6) {$PQ^2=(x_2-x_1)^2+(y_2-y_1)^2$};
\end{tikzpicture}\end{center}

当 $B\ne0$ 时，$k=-A/B$，有 $y_2-y_1=k(x_2-x_1)$。由勾股定理，
\[
 PQ=\sqrt{1+k^2}\,|x_2-x_1|.
\]
根差不必求出两个根：
\[
 (x_2-x_1)^2=(x_1+x_2)^2-4x_1x_2=\frac{\Delta_x}{D^2}.
\]
所以 $PQ=\sqrt{1+k^2}\sqrt{\Delta_x}/D$，代入 $\Delta_x=4\alpha\beta B^2H$，
与无斜率的通用弦长公式一致。若方程整体约去公因数，必须同时使用\textbf{新系数及其判别式}；不能混用。
\end{proofbox}
